\subsection{Positional Systems}

\content{
        As we have noticed, the notion of `position' is extremely valuable, but
        it took a long time for this idea to develop. Let's discuss some of the
        basics of what a posisitional system is. 
}{% presentation
\vspace{3in}
}{% nopresentation
\vspace{3in}
}{% comments
Write a number, explain breifly what each position is, this is pretty basic.
Note that ancient civilizations did not neccesarily use 10 as their base. 
}

\content{
        If we choose 8 as our base for a posistional system, what are the values
        of each position?
}{% presentation
\vspace{2in}
}{% nopresentation
\vspace{2in}
}{% comments
$8^0=1, 8^1=8, 8^2=64, 8^3=512, 8^4=4096,\dots$. 
}

\content{
        Write the number $502$ (which is in base 8) in base 10. 
}{% presentation
\vspace{2in}
}{% nopresentation
\vspace{2in}
}{% comments
make a notational note that we will always write non-base 10 numbers with the
base as a subscript. So this number should be written $502_8$.
}

\content{
        Write the number $3324_5$ in base 10. 
}{% presentation
\vspace{3in}
}{% nopresentation
\vspace{3in}
}{% comments
}

\content{
        Write the number $10101_2$ in base 10. 
}{% presentation
\vspace{3in}
}{% nopresentation
\vspace{3in}
}{% comments
}

